\section{Introduction}
Space manipulators are central to on-orbit servicing, debris removal, and maintenance of failed spacecraft. Unlike ground-based manipulators, a pre-capture space manipulator is free floating: joint motion induces base reaction motion through momentum conservation and changes the end-effector trajectory. Its dynamics are nonlinear and strongly coupled, with parameter uncertainty, time-varying disturbances, and actuator faults. Accurate tracking with a prescribed convergence time is therefore a fundamental problem.

\subsection{Related Work and Background}
Under zero initial linear and angular momentum, base velocity is determined by joint velocity, and end-effector velocity is described by the generalized Jacobian rather than a fixed-base Jacobian. The generalized Jacobian was introduced by Umetani and Yoshida~\cite{umetani1989}; the virtual manipulator approach characterized the coupled kinematics and dynamics~\cite{vafa1990}. Dynamic singularities and the properties of free-floating and free-flying robots were studied in~\cite{papadopoulos1993,dubowsky1993}. Adaptive computed-torque and reactionless-motion controllers were developed in~\cite{parlaktuna2004,xu2016}. In particular, the reduced dynamics can be written as
\[
H_g(q)\ddot q+C_g(q,\dot q)\dot q=\tau_m.
\]
Trajectory planning, model predictive control, time-delay estimation, adaptive sliding-mode control, and prescribed-performance control have also been investigated~\cite{wang2015,ning2019,yu2022,yao2021}. Their convergence times, however, generally depend on gains and initial conditions.

Fixed-time and prescribed-time methods provide convergence bounds selected at design time~\cite{polyakov2012,zuo2019}. Robot tracking and space-robot pose control have been addressed using these ideas~\cite{su2020,xuPredefinedTimeTrackingControl2024,li2025}, but nonsmooth terminal surfaces or fractional powers can cause input jumps or gain singularities. Pyragas introduced delayed feedback control~\cite{pyragas1992}; smooth periodic delayed feedback (PDF) provides fixed-time stabilization for controllable linear time-delay systems while avoiding terminal gain singularities~\cite{zhouFixedtimeStabilizationLinear2022}. Its mechanism is the annihilation or nilpotence of a finite-dimensional one-period state-transition matrix. Its use for nonlinear manipulators requires accurate disturbance compensation. Prescribed-time disturbance observers provide an appropriate interface~\cite{jiangPrescribedtimeDisturbanceObserver2024}.

\subsection{Contributions and Organization}
This paper develops a four-stage structure: disturbance observation, dynamic compensation, error linearization, and periodic delayed feedback. The contributions are: (1) a formal reduction from base reaction motion and momentum conservation to equivalent joint-space dynamics; (2) a prescribed-time disturbance observer and a two-stage timing design; (3) a normalized controllable error channel for PDF design; (4) a Gramian-based gain construction and explicit practical error bounds; and (5) MATLAB validation with a seven-degree-of-freedom model.

The remainder is organized as follows. Section~\ref{sec:modeling} establishes the model; Section~\ref{sec:problem} formulates the problem; Section~\ref{sec:controller} designs the observer, compensation, and PDF controller; Section~\ref{sec:stability} proves convergence; Section~\ref{sec:simulation} presents simulations; and Section~\ref{sec:conclusion} concludes the paper.
